\documentclass[11pt]{article}
\usepackage[margin=0.88in]{geometry}
\usepackage{amsmath,amssymb,lmodern}
\usepackage[T1]{fontenc}
\usepackage{needspace}
\usepackage[colorlinks=true,linkcolor=blue,urlcolor=blue]{hyperref}
\setlength{\emergencystretch}{3em}
\setlength{\parskip}{0.35em}
\setcounter{secnumdepth}{0}
\providecommand{\tightlist}{\setlength{\itemsep}{0pt}\setlength{\parskip}{0pt}}
\title{Sharp total variation bounds from spherical rearrangement}
\author{Henry Zweiman}
\date{September 15, 2026}
\begin{document}
\maketitle
\Needspace{8\baselineskip}\section{Abstract}\label{abstract}

We study the sharp control of total variation by a higher derivative
under homogeneous endpoint conditions. For every integer \(n\ge2\), we
prove that the norm of \(u\mapsto u'\) from \(W_0^{n,2}(0,1)\) to
\(L^1(0,1)\) is twice the norm of point evaluation into
\(L^\infty(0,1)\). The extremizers of the two embeddings coincide and
are symmetric about the midpoint. This proves the Hilbert-space case of
a conjecture of Nazarov and Shcheglova. The main step is a rearrangement
principle for a clamped Green operator after projection onto the
zero-mean subspace. A Jacobi-polynomial representation identifies its
quadratic form with a positive mixture of Poisson interactions on a
sphere, where rearrangement reduces the optimization to a cap. We also
prove the constant identity and coincidence of maximizing sets at the
measure endpoint \(p=1\), using the geometry of extreme moment measures
and unimodal splines. For each fixed derivative order, we further prove
the full variation identity and midpoint-symmetric equality
classification on an open interval of exponents around two. A dual
optimality equation gives compactness through the highest derivative;
polynomial sublevel bounds control the point-evaluation profile on both
sides of the Hilbert exponent. Symmetry at the measure endpoint and the
full all-exponent conjecture remain open.

\Needspace{8\baselineskip}\section{1. The sharp inequality}\label{the-sharp-inequality}

All functions and measures in this paper are real. For
\(1\le p\le\infty\), we use the homogeneous norm \(\|u^{(n)}\|_p\) on
\(W_0^{n,p}(0,1)\). The endpoint conditions are

\[
u^{(j)}(0)=u^{(j)}(1)=0,\qquad 0\le j<n.
\tag{1}
\]

Let \(C_{n,p}\) and \(V_{n,p}\) denote the sharp constants in

\[
\|u\|_\infty\le C_{n,p}\|u^{(n)}\|_p,
\qquad
\|u'\|_1\le V_{n,p}\|u^{(n)}\|_p.
\tag{2}
\]

The inequality \(V_{n,p}\ge2C_{n,p}\) follows from the variation of a
function that vanishes at both endpoints. Nazarov and Shcheglova {[}NS,
Conjecture 4.14{]} conjecture equality for all \(n\ge2\) and
\(1\le p\le\infty\), together with coincidence and midpoint symmetry of
the extremizers. Their Remark 4.15 records the previously known case
\(n=2\).

\Needspace{6\baselineskip}\textbf{Theorem 1.} For every integer \(n\ge2\),

\[
V_{n,2}=2C_{n,2}
=\frac{1}{2^{2n-2}(n-1)!\sqrt{2n-1}}.
\tag{3}
\]

The nonzero extremizers are precisely the scalar multiples of the Riesz
representer \(R_n(\cdot,1/2)\) of evaluation at \(1/2\) in
\(W_0^{n,2}(0,1)\). In particular, the maximizing sets in (2) coincide,
and their elements are even about \(1/2\).

The point-evaluation formula used here is prior work. Kalyabin {[}K,
Theorem 2{]} gives, after affine rescaling,

\[
\|\operatorname{ev}_a\|^2
=\frac{[a(1-a)]^{2n-1}}{(2n-1)((n-1)!)^2}.
\tag{4}
\]

Our contribution is the upper bound for the variation norm and its
equality classification, obtained through the rearrangement principle in
Section 3. Best polynomial approximation already occurs in this
embedding problem {[}GS{]}. The sharp \(W_0^{m,2}\) to \(L^1\) embedding
of the function itself was proved by Hindov, Nitzan, Olsen and Rydhe
{[}HNOR{]}; their theorem does not include the additional zero-mean
constraint on a derivative that appears below.

\Needspace{8\baselineskip}\section{2. A Jacobi representation of the clamped Green
form}\label{a-jacobi-representation-of-the-clamped-green-form}

Fix \(m\ge1\), and set

\[
w(x)=x^m(1-x)^m,\qquad
M=\int_0^1w(x)\,dx=\frac{(m!)^2}{(2m+1)!},
\qquad d\pi(x)=M^{-1}w(x)\,dx.
\tag{5}
\]

Let \(h_k\) be the orthonormal polynomials for \(\pi\), with positive
leading coefficients. These are normalized shifted Jacobi polynomials
with parameters \((m,m)\); in particular \(h_0=1\). Put

\[
v_k=wh_k,\qquad \lambda_k=\frac{(k+2m)!}{k!}.
\tag{6}
\]

\Needspace{6\baselineskip}\textbf{Lemma 2.} The functions \(v_k/\sqrt{M\lambda_k}\) form an
orthonormal basis of \(H=W_0^{m,2}(0,1)\) in its energy inner product.
Moreover,

\[
(-1)^m v_k^{(2m)}=\lambda_k h_k.
\tag{7}
\]

\emph{Proof.} The left side of (7) has degree \(k\). If \(q\) is a
polynomial of degree less than \(k\), integration by parts \(2m\) times
yields

\[
\int_0^1 (wh_k)^{(2m)}wq
=\int_0^1 wh_k(wq)^{(2m)}=0.
\tag{8}
\]

Every boundary term vanishes because one of its two factors has fewer
than \(m\) derivatives and retains an endpoint zero. Orthogonality
therefore shows that the left side of (7) is proportional to \(h_k\).
Comparing leading coefficients proves (7). A further integration by
parts gives

\[
\langle v_k,v_\ell\rangle_H=M\lambda_k\delta_{k\ell}.
\tag{9}
\]

If \(y\in H\) is orthogonal to every \(v_k\), weak integration by parts
gives \(\int_0^1 yh_k=0\) for every \(k\). The \(h_k\) span the
polynomials, which are dense in \(C[0,1]\), so \(y=0\). This proves
completeness. \(\square\)

The mean functional \(y\mapsto\int_0^1y\) has Riesz representer
\(v_0/\lambda_0=w/(2m)!\). Therefore

\[
H_*:=\left\{y\in H:\int_0^1y=0\right\}
\tag{10}
\]

is obtained by removing the zeroth basis vector. For a bounded real load
\(g\), Parseval gives the squared dual norm of \(y\mapsto\int gy\) on
\(H_*\):

\[
Q_m(g)
=M\sum_{k\ge1}\frac{\left|\int_0^1gh_k\,d\pi\right|^2}{\lambda_k}.
\tag{11}
\]

Equivalently, if \(G_m\) is the clamped Green kernel and
\(w_m=w/(2m)!\), then

\[
Q_m(g)=\iint g(x)G_m(x,t)g(t)\,dx\,dt
-\frac{\left(\int gw_m\right)^2}{\int w_m}.
\tag{12}
\]

Formula (11) is sufficient for the proof and avoids any pointwise
convergence issue for a Green series.

\Needspace{8\baselineskip}\section{3. Spherical rearrangement of the
load}\label{spherical-rearrangement-of-the-load}

Set \(d=2m+2\) and let \(\sigma\) be normalized surface measure on
\(S^d\). The coordinate \(x=(1+\xi_{d+1})/2\) has law \(\pi\). Thus

\[
F(\xi)=g\!\left(\frac{1+\xi_{d+1}}2\right)
\tag{13}
\]

identifies \(L^2(\pi)\) with the zonal subspace of \(L^2(S^d)\). Under
this identification, \(h_k\) is the normalized zonal harmonic of degree
\(k\). This follows from the orthogonal decomposition of polynomials
into spherical harmonics: a zonal polynomial of degree \(k\) orthogonal
to all lower zonal degrees has only its degree-\(k\) component.

The spherical Poisson operator \(P_r\), \(0<r<1\), has eigenvalue
\(r^k\) on degree-\(k\) harmonics and kernel

\[
K_r(\xi\cdot\eta)
=\frac{1-r^2}{(1-2r\xi\cdot\eta+r^2)^{(d+1)/2}}.
\tag{14}
\]

We use normalized surface measure, as in {[}ABR, (1.15){]}. The
eigenvalue follows by harmonic extension of a homogeneous harmonic
polynomial. The beta integral gives

\[
\frac1{\lambda_k}
=\frac1{(2m-1)!}\int_0^1r^k(1-r)^{2m-1}\,dr.
\tag{15}
\]

Tonelli\textquotesingle s theorem applied to the nonnegative spectral
terms in (11) therefore gives

\[
Q_m(g)=\frac{M}{(2m-1)!}\int_0^1(1-r)^{2m-1}
\left(\langle F,P_rF\rangle-\left(\int F\,d\sigma\right)^2\right)dr.
\tag{16}
\]

The deletion of the constant harmonic is exactly the mean subtraction
needed in (12).

\Needspace{6\baselineskip}\textbf{Theorem 3.} If \(A\subset(0,1)\) is measurable and \(a\in[0,1]\)
satisfies \(\pi(A)=\pi([0,a])\), then

\[
Q_m(\mathbf1_A)\le Q_m(\mathbf1_{[0,a]}).
\tag{17}
\]

If \(0<\pi(A)<1\), equality holds only when \(A\) agrees, up to null
sets, with an interval abutting one endpoint and having that
\(\pi\)-measure.

\emph{Proof.} Lift \(A\) to the zonal subset \(E\subset S^d\) using
(13). For fixed \(r\), \(K_r\) is positive, bounded and strictly
decreasing in geodesic distance. Spherical rearrangement {[}BH, Theorem
2{]} maximizes the self-interaction of \(\mathbf1_E\) at fixed surface
measure by a cap. Choose that cap about the negative coordinate axis. It
corresponds to \([0,a]\). The mean term in (16) is unchanged, and
integration proves (17).

For clarity, the equality hypotheses in {[}BH{]} are satisfied with two
functions and integrand \((s,t)\mapsto st\). Its mixed finite difference
is strictly positive; the kernel is positive and strictly decreasing;
and the nontrivial indicators are nonconstant with finite interaction.
Hence equality at fixed \(r\) forces \(E\) to be a cap up to rotation.
If equality holds in (17), the nonnegative rearrangement deficit has
zero integral in (16), so equality holds for almost every \(r\), and in
particular for one \(r\in(0,1)\). Thus \(E\) is a cap.

Its axis must be the original coordinate axis up to reversal. Indeed,
zonal invariance makes the centroid of \(E\) parallel to that axis. A
nontrivial cap has a nonzero centroid parallel to its own axis.
Consequently \(A\) is an endpoint interval, as claimed. \(\square\)

The same argument proves a broader load rearrangement statement. For a
bounded real \(g\), let \(g^\downarrow\) be its nonincreasing
equimeasurable rearrangement relative to \(\pi\). Adding a constant
makes \(g\) nonnegative without changing \(Q_m\). Applying {[}BH{]} in
(16) gives

\[
Q_m(g)\le Q_m(g^\downarrow).
\tag{18}
\]

This is rearrangement of a load in a Green quadratic form. It does not
assert that ordinary rearrangement preserves a higher derivative norm.

\Needspace{8\baselineskip}\section{4. Proof of Theorem 1}\label{proof-of-theorem-1}

Put \(m=n-1\). The map \(u\mapsto y=u'\) is an isometric bijection from
\(W_0^{n,2}\) to \(H_*\): integrating \(y\) recovers \(u\), and the
zero-mean condition supplies \(u(1)=0\).

Duality gives

\[
V_{n,2}^2=\sup_{|g|\le1}Q_m(g)
=4\sup_A Q_m(\mathbf1_A).
\tag{19}
\]

To justify the second equality, write a bounded \(g\) with \(|g|\le1\)
as the average of \(\operatorname{sign}(g-s)\) over uniform
\(s\in[-1,1]\). Convexity of the squared Hilbert dual norm bounds its
value by the supremum over sign functions. Constants are annihilated in
(11), giving \(Q_m(2\mathbf1_A-1)=4Q_m(\mathbf1_A)\).

For an endpoint interval,

\[
\int_0^a y(x)\,dx=u(a),
\qquad Q_m(\mathbf1_{[0,a]})=\|\operatorname{ev}_a\|^2.
\tag{20}
\]

Theorem 3, (4), and (19) now prove (3). In particular, the maximizing
interval is a half-interval.

Suppose that \(u\ne0\) attains equality, and take
\(g=\operatorname{sign}(u')\), choosing either sign where \(u'=0\). Both
signs occur on sets of positive measure. Equality in the Hilbert dual
bound and in Theorem 3 forces \(g\) to be a half-interval sign function,
up to its global sign and null sets. Its functional on \(H_*\) is
\(\pm2u(1/2)\). Equality in Cauchy-\/-Schwarz therefore makes \(u\) a
scalar multiple of \(R_n(\cdot,1/2)\).

Conversely, this representer maximizes point evaluation. Its variation
is at least twice its maximum and at most the bound just proved, so it
also maximizes variation. Formula (4) has its unique maximum at
\(a=1/2\), and uniqueness of a Hilbert-space Riesz representer gives
precisely the same classification for the point-norm extremizers.
Finally, reflection preserves the energy inner product and fixes
evaluation at \(1/2\), so uniqueness makes \(R_n(\cdot,1/2)\) reflection
invariant. \(\square\)

\Needspace{8\baselineskip}\section{5. The measure endpoint}\label{the-measure-endpoint}

Let \(\mathcal M_n\) consist of real finite signed measures \(\mu\) on
\([0,1]\) such that

\[
\|\mu\|_{\mathrm{TV}}\le1,
\qquad \int_0^1t^j\,d\mu(t)=0\quad(0\le j<n).
\tag{21}
\]

Their compactly supported potentials are

\[
U_\mu(x)=\int_0^1\frac{(x-t)_+^{n-1}}{(n-1)!}\,d\mu(t).
\tag{22}
\]

The moments make \(U_\mu\) vanish off \([0,1]\), and \(D^nU_\mu=\mu\).
For \(n\ge2\), \(U_\mu\) is continuous and absolutely continuous, with
integrable first derivative. Endpoint atoms are allowed in this relaxed
formulation.

\Needspace{6\baselineskip}\textbf{Theorem 4.} For every \(n\ge2\),

\[
\sup_{\mu\in\mathcal M_n}\|U_\mu'\|_1
=2\sup_{\mu\in\mathcal M_n}\|U_\mu\|_\infty.
\tag{23}
\]

The maximizing measures for the two quantities coincide. The relaxed
constants equal \(V_{n,1}\) and \(C_{n,1}\), respectively. No symmetry
assertion at \(p=1\) is included.

\emph{Proof of the constant identity.} Every extreme point of
\(\mathcal M_n\) has norm one and exactly \(n+1\) support points. To
prove the support bound, suppose \(\mu\) has norm one and at least
\(n+2\) disjoint Borel sets of positive \(|\mu|\)-measure. A nonzero
function \(h\) constant on those sets can satisfy the \(n\) equations
\(\int t^jh\,d\mu=0\) and the extra equation \(\int h\,d|\mu|=0\). Then
\((1\pm\varepsilon h)\mu\) are distinct norm-one admissible measures for
small \(\varepsilon\). Such a \(\mu\) is not extreme. A nonzero measure
of norm less than one is excluded by scalar perturbation. Zero is the
midpoint of opposite nonzero moment measures, which exist on any \(n+1\)
distinct nodes, and is not extreme either. A nonzero measure on at most
\(n\) points is excluded by the Vandermonde determinant.

On nodes \(t_0<\cdots<t_n\), the moment nullspace is one-dimensional,
with weights

\[
b_i=\prod_{j\ne i}(t_i-t_j)^{-1}.
\tag{24}
\]

Thus every extreme measure is a normalized signed multiple of
\(\sum_i b_i\delta_{t_i}\). Its potential is a B-spline, up to scale and
sign. More explicitly, the Genocchi-\/-Hermite identity {[}dB, (52){]}
identifies

\[
\rho(x)=n\sum_{i=0}^n b_i(t_i-x)_+^{n-1}
=(-1)^n n!\,U_{\sum b_i\delta_{t_i}}(x)
\tag{25}
\]

with the density of \(\sum_i t_iX_i\), where \(X\) is uniform on the
standard \(n\)-simplex. After parallel sections of the simplex are
identified by translation, their Minkowski interpolation is contained in
the section at the interpolated coordinate. The classical
Brunn-\/-Minkowski inequality {[}GHW, (3), \(p=1\){]} therefore makes
their volumes to the power \(1/(n-1)\) concave. The coarea factor
relating section volume to \(\rho\) is constant, so \(\rho^{1/(n-1)}\)
is concave and \(\rho\) is unimodal. It is positive between its end
knots, continuous, and zero at those knots. Every extreme potential
therefore satisfies

\[
\|U_\mu'\|_1=2\|U_\mu\|_\infty.
\tag{26}
\]

Write \(C=\sup_{\mathcal M_n}\|U_\mu\|_\infty\). The set
\(\mathcal M_n\) is weak-star compact. Moreover,
\(\mu\mapsto\|U_\mu'\|_1\) is weak-star lower semicontinuous: it is the
supremum of the continuous linear functionals \(\int H_g\,d\mu\), where
\(|g|\le1\) and

\[
H_g(t)=\int_t^1\frac{(x-t)^{n-2}}{(n-2)!}g(x)\,dx.
\tag{27}
\]

Its sublevel set at \(2C\) is closed and convex and contains all extreme
points by (26). Krein-\/-Milman therefore gives the upper bound in (23).
The elementary variation lower bound gives the reverse inequality. The
height supremum is attained: on the bounded measure set, joint
continuity of the kernels in (22) makes the potential map weak-star to
uniform continuous. Consequently the variation supremum is attained as
well.

\emph{Coincidence.} Let \(\mu\) maximize variation, choose
\(g=\operatorname{sign}(U_\mu')\), and let \(\mathcal F\) be the compact
exposed face on which \(\int H_g\,d\nu=2C\). Its extreme points are
extreme in \(\mathcal M_n\). For each such point \(e\),

\[
2C=\int gU_e'\le\|U_e'\|_1=2\|U_e\|_\infty\le2C.
\tag{28}
\]

Thus \(g\) agrees with the sign of \(U_e'\) wherever the latter is
nonzero, and every \(U_e\) maximizes height. Its end knots must be \(0\)
and \(1\): a potential supported on a shorter interval of length \(L\)
can be stretched to \([0,1]\), keeping its derivative-measure norm one
and multiplying its height by \(L^{1-n}>1\).

All \(U_e\) have the same global sign, since opposite signs give
opposite nonzero derivatives on a common interval near zero,
contradicting their agreement with \(g\). Change that common sign to
positive. Concavity of \(U_e^{1/(n-1)}\) makes its maximum set a closed
interval, with strictly positive derivative almost everywhere before it
and strictly negative derivative almost everywhere after it. Two
disjoint maximum intervals would again force contradictory signs for
\(g\). The maximum intervals therefore intersect pairwise, and hence
have a common point \(a_*\). All \(U_e(a_*)=C\). Evaluation is weak-star
continuous, so Krein-\/-Milman applied to \(\mathcal F\) gives
\(U_\mu(a_*)=C\). Conversely, a height maximizer is a variation
maximizer by the already proved constant identity.

\emph{Recovery of the ordinary constants.} For an ordinary admissible
\(u\), the measure \(u^{(n)}dx\) satisfies (21). Conversely, mollify the
zero extension of \(U_\mu\) with a nonnegative mollifier supported in
\([-\varepsilon,\varepsilon]\), obtaining \(U_\varepsilon\). With
\(a=1+2\varepsilon\), set

\[
u_\varepsilon(x)=a^{1-n}U_\varepsilon(ax-\varepsilon).
\tag{29}
\]

It is smooth, supported on \([0,1]\), satisfies all endpoint conditions,
and has \(\|u_\varepsilon^{(n)}\|_1\le1\). Its maximum and variation
converge to those of \(U_\mu\), by uniform convergence of the potential
mollification and \(L^1\) convergence of its first derivative. This
proves equality of the relaxed and ordinary sharp constants. The proof
classifies coincidence of relaxed maximizers and does not assert
attainment in the ordinary \(W_0^{n,1}\) class. \(\square\)

\Needspace{8\baselineskip}\section{6. An open interval of
exponents}\label{an-open-interval-of-exponents}

The Hilbert theorem also yields a neighborhood of exponents for each
fixed derivative order. The interval is not asserted to be uniform in
the order. In this section write \(X_{n,p}=W_0^{n,p}(0,1)\) for the
clamped space defined by (1), with norm \(\|u^{(n)}\|_p\).

\Needspace{6\baselineskip}\textbf{Theorem 5 (local continuation in the exponent).} For every fixed
integer \(n\ge3\), there exists \(0<\varepsilon_n<1\) such that, if
\(|p-2|<\varepsilon_n\), then

\[
V_{n,p}=2C_{n,p}.
\tag{6.1}
\]

All nonzero extremizers in the two inequalities are the scalar multiples
of the unique positive norm-one extremizer for evaluation at \(1/2\).
They are symmetric about the midpoint and, after a choice of global
sign, have strictly positive derivative on \((0,1/2)\) and strictly
negative derivative on \((1/2,1)\).

The radius is allowed to depend on \(n\) and is not estimated here. The
theorem does not reach all \(1<p<\infty\), give an interval uniform in
\(n\), or settle the remaining endpoint symmetry questions.

\subsection{6.1. Dual representation and compactness of variation
maximizers}\label{dual-representation-and-compactness-of-variation-maximizers}

Let \(\mathcal P_{n-1}\) denote the polynomials of degree at most
\(n-1\). Integration from the left endpoint identifies \(X_{n,p}\)
isometrically with

\[
Z_{n,p}=\{f\in L^p(0,1):\int_0^1t^jf(t)\,dt=0,\ 0\le j<n\}.
\tag{6.2}
\]

Indeed \(u(x)=\int_0^x(x-t)^{n-1}f(t)/(n-1)!\,dt\); the right endpoint
conditions are equivalent to the displayed moments. For bounded \(g\),
put

\[
H_g(t)=\int_t^1\frac{(x-t)^{n-2}}{(n-2)!}g(x)\,dx.
\tag{6.3}
\]

Fubini gives \(\int gu'=\int H_gf\). If \(q=p/(p-1)\), the norm of this
functional on \(Z_{n,p}\) equals

\[
\min_{P\in\mathcal P_{n-1}}\|H_g-P\|_q.
\tag{6.4}
\]

Here is a direct justification without a quotient-duality assumption. A
minimizing polynomial exists by finite-dimensional coercivity and is
unique for \(q>1\). Its residual \(R=H_g-P\) satisfies

\[
\int_0^1 |R|^{q-2}R\,Q=0\qquad(Q\in\mathcal P_{n-1}).
\tag{6.5}
\]

If \(R\ne0\), the function \(f=\operatorname{sign}(R)|R/\|R\|_q|^{q-1}\)
has \(L^p\) norm one and belongs to (6.2). It realizes the bound in
Hölder\textquotesingle s inequality. The zero-residual case is
immediate. This proves (6.4).

For every fixed \(p>1\), the variation supremum is attained. A
norm-bounded sequence is bounded in \(W^{n,p}\) by successive
integration, and its derivatives through order \(n-1\) are uniformly
bounded and equicontinuous, using Hölder\textquotesingle s inequality
for the highest integral. A subsequence converges uniformly through
order \(n-1\), and the highest derivatives have a weakly convergent
subsequence in \(L^p\). The limit satisfies all endpoint conditions, has
norm at most one, and retains the limiting variation. The supremum is
positive, so a maximizer must have norm exactly one by scaling.

For a norm-one variation maximizer \(u\), choose
\(g=\operatorname{sign}(u')\), with any value in \([-1,1]\) on the zero
set. The functional in (6.4) is at most \(V_{n,p}\), since
\(\int gv'\le\|v'\|_1\). On \(u\) it equals \(V_{n,p}\). Thus its
residual satisfies \(\|R\|_q=V_{n,p}\), and equality in Hölder gives

\[
u^{(n)}(t)=\operatorname{sign}R(t)
\left|\frac{R(t)}{V_{n,p}}\right|^{q-1}.
\tag{6.6}
\]

This formula gives more compactness than the original Sobolev bound.

\Needspace{6\baselineskip}\textbf{Lemma 6 (compactness of maximizers).} Suppose \(p_j\to2\) and
\(u_j\) is any sequence of norm-one variation maximizers in
\(X_{n,p_j}\). Then, after a subsequence and a choice of global signs,
\(u_j\) converges in \(C^n[0,1]\) to the positive norm-one Hilbert
maximizer \(u_*\).

\emph{Proof.} Restrict to \(3/2\le p_j\le5/2\), so \(5/3\le q_j\le3\).
The functions \(H_{g_j}\), with \(|g_j|\le1\), are uniformly bounded and
Lipschitz. For \(n=2\), this follows directly from \(H_g'=-g\); for
\(n\ge3\), differentiate (6.3) once and bound its remaining integral.

The minimizing polynomials satisfy
\(\|P_j\|_{q_j}\le2\|H_{g_j}\|_{q_j}\). Their coefficients are uniformly
bounded. To verify the needed uniform norm equivalence, a contrary
sequence of polynomials can be normalized to have coefficient norm one
and \(L^{q_j}\) norm tending to zero. A coefficient-convergent
subsequence would give a nonzero polynomial vanishing almost everywhere,
a contradiction. Consequently both \(P_j\) and their derivatives are
uniformly bounded, and the residuals \(R_j\) are uniformly bounded and
Lipschitz.

The quantities \(V_{n,p_j}\) are bounded above by (6.4) with the zero
polynomial. They are bounded away from zero by testing the one fixed
nonzero function \(u_*\), whose highest derivative is bounded. Formula
(6.6) now makes \(u_j^{(n)}\) uniformly bounded and equicontinuous. For
example, signed power maps with exponent in \([2/3,2]\) on a fixed
bounded interval have a common Hölder modulus with exponent \(2/3\),
after enlarging the constant. Arzelà-\/-Ascoli supplies a uniformly
convergent subsequence of the highest derivatives. Integration and the
left endpoint conditions give convergence in \(C^n\) to some \(u\).

Uniform boundedness and \(p_j\to2\) imply
\(\int |u^{(n)}|^2=\lim_j\int|u_j^{(n)}|^{p_j}=1\). All right endpoint
conditions pass to the limit. Also \(\|u'\|_1=\lim_jV_{n,p_j}\). Testing
\(u_*/\|u_*^{(n)}\|_{p_j}\) gives \(\liminf_j V_{n,p_j}\ge V_{n,2}\),
while the Hilbert inequality applied to the limit gives the reverse
bound. Thus \(u\) is a Hilbert maximizer. Theorem 1 identifies it with
\(u_*\) or \(-u_*\), proving the lemma. \(\square\)

\subsection{6.2. Persistence of a single
peak}\label{persistence-of-a-single-peak}

We give the details that weak or uniform convergence alone would miss.
The Hilbert midpoint representer, up to a positive factor, can be
written

\[
U(x)=F(|1-2x|),\qquad
F(z)=\int_z^1(t^2-z^2)^{n-1}\,dt,
\quad 0\le z\le1.
\tag{6.7}
\]

To check this representation, expand the integrand. The resulting
expression for \(F\) is an even polynomial of degree at most \(2n-2\),
plus a multiple of \(z^{2n-1}\). Thus (6.7) is polynomial of degree at
most \(2n-1\) on each half-interval and has matching derivatives through
order \(2n-2\) at the midpoint. Near \(z=1\), writing \(t=z+(1-z)s\)
shows that \(F(z)\) is a positive constant times \((1-z)^n\), to leading
order. Therefore all endpoint derivatives below order \(n\) vanish. The
jump of the derivative of order \(2n-1\) is nonzero. The weak derivative
\((-1)^nD^{2n}U\) is a scalar multiple of the midpoint delta. Testing
against \(U\) shows this scalar is positive because
\(\int|U^{(n)}|^2>0\) and \(U(1/2)>0\). Uniqueness of the Riesz
representer proves the claim up to positive normalization.

For \(n\ge2\), differentiation gives

\[
F'(z)=-2(n-1)z\int_z^1(t^2-z^2)^{n-2}\,dt<0
\qquad(0<z<1).
\tag{6.8}
\]

Also \(F''(0)=-2(n-1)/(2n-3)<0\). Hence \(u_*'>0\) on \((0,1/2)\),
\(u_*'<0\) on \((1/2,1)\), and \(u_*''(1/2)<0\). The endpoint expansion
gives \(u_*^{(n)}(0)>0\) and \((-1)^nu_*^{(n)}(1)>0\).

These properties imply that every sufficiently close \(C^n\)
perturbation satisfying the same clamped endpoint conditions has exactly
one interior critical point, a strict maximum. Indeed, near zero its
highest derivative remains positive and

\[
u'(x)=\int_0^x\frac{(x-t)^{n-2}}{(n-2)!}u^{(n)}(t)\,dt>0.
\tag{6.9}
\]

The reflected argument gives negativity near one. On compact intervals
away from the endpoints and midpoint the strict derivative signs
persist. On a small interval about the midpoint, the second derivative
stays negative, so there is precisely one zero of the first derivative,
bracketed by its signs on either side.

Lemma 6 therefore implies that for all \(p\) sufficiently close to two,
every norm-one variation maximizer is, up to sign, positive and strictly
unimodal. Such a maximizer has variation exactly twice its height. Since
every admissible function has variation at least twice its height, this
proves (6.1). It also proves coincidence of maximizing sets: a variation
maximizer has height \(C_{n,p}\), and every height maximizer has
variation between \(2C_{n,p}\) and \(V_{n,p}\).

The only conclusion still needed for the local theorem is midpoint
symmetry. Single-peak geometry by itself does not supply it.

\subsection{6.3. Smoothness of the point-evaluation profile near
two}\label{smoothness-of-the-point-evaluation-profile-near-two}

Let \(m=n-1\ge2\), and define

\[
K_a(t)=\frac{(a-t)_+^m}{m!},\qquad
J(q,a)=\min_{P\in\mathcal P_m}\int_0^1|K_a-P|^q\,dt.
\tag{6.10}
\]

By the same moment argument as (6.4), the evaluation norm is
\(J(q,a)^{1/q}\). We need uniform continuity on the whole interval in
\(a\), and twice continuous differentiability in \(a\) near
\((q,a)=(2,1/2)\).

First, minimizing polynomials have coefficients bounded uniformly for
\(q\) in a compact subset of \((1,\infty)\) and \(a\in[0,1]\), by the
argument used above. On this bounded coefficient set the integrands in
(6.10) vary continuously with \((q,a,P)\), uniformly in \(t\).
Minimizing over that compact set proves continuity of \(J\), including
uniform convergence \(J(q,\cdot)\to J(2,\cdot)\) as \(q\to2\). Strict
convexity gives a unique minimizing polynomial; compactness and
uniqueness prove continuous dependence of its coefficients.

The following elementary sublevel estimate handles exponents just below
two. If \(P\) is a polynomial of degree at most \(m\) on an interval
\(I\) of length bounded above and below, and
\(\|P\|_{L^\infty(I)}\ge b>0\), then

\[
|\{t\in I:|P(t)|\le h\}|\le C(m,b,I)h^{1/m}
\qquad(0<h<1).
\tag{6.11}
\]

The constant is uniform for intervals in the stated length range. To
prove this, let the sublevel set have length \(\ell>0\). Divide its mass
into sufficiently many equal portions and choose \(m+1\) points from
alternate portions. They can be chosen in the sublevel set with
consecutive separations bounded below by a fixed multiple of
\(\ell/(m+1)\): coordinate separation is at least the measure of the
intervening portion. Lagrange interpolation at these points gives
\(\|P\|_\infty\le C_m h\ell^{-m}\), with the interval-length factor
absorbed into the constant. Rearrangement gives (6.11); approximation of
the chosen points handles essential rather than attained sublevel
membership.

Consequently, for \(0<s<1/m\), integration of distribution functions
gives

\[
\int_{\{|P|<h\}}|P|^{-s}\,dt\le C h^{1/m-s}.
\tag{6.12}
\]

For example, split the region into dyadic bands
\(2^{-k-1}h\le|P|<2^{-k}h\), apply (6.11) to each, and sum the
convergent geometric series. The estimates also show uniform
integrability of these inverse powers.

At \((q,a)=(2,1/2)\), the minimizing residual in (6.10) is a nonzero
polynomial on each of the two half-intervals. If it were identically
zero on one half, its orthogonality to constants would fail on the
other: the remaining residual would be a nonzero constant multiple of
the nonnegative truncated power. Continuity of the coefficients
therefore supplies a uniform positive lower bound for the sup norm of
each polynomial piece on its own interval when \((q,a)\) and the
polynomial coefficients are sufficiently close to this point. The
intervals\textquotesingle{} lengths also stay bounded away from zero.
Thus (6.11)-\/-(6.12) apply uniformly to both pieces.

Write \(P_c(t)=\sum_{j=0}^mc_jt^j\) and \(R=K_a-P_c\). The normal
equations are

\[
\mathcal N_j(q,a,c)=\int_0^1|R|^{q-2}R\,t^j\,dt=0,
\qquad 0\le j\le m.
\tag{6.13}
\]

Take a small open neighborhood with \(q>2-1/(2m)\). The equations are
continuously differentiable in \((q,a,c)\) there. Their derivatives in
\((a,c)\) are obtained by multiplying the variation of \(R\) by
\((q-1)|R|^{q-2}\), and the derivative in \(q\) has integrand
\(|R|^{q-2}R\log|R|\,t^j\), with value zero at \(R=0\). These assertions
follow first away from a small sublevel set of \(R\), by ordinary
differentiation. The uniform bound (6.12) makes the omitted derivatives
uniformly integrable and arbitrarily small as the sublevel threshold
tends to zero. The logarithmic derivative is bounded by a constant times
a small positive power of \(|R|\) near zero. This also proves continuity
of the derivatives by splitting into the same two regions. The moving
interface \(t=a\) causes no additional term, since \(K_a\) and its first
\(a\)-derivative are continuous there for \(m\ge2\).

At the base point, the coefficient derivative of (6.13) is the negative
definite Gram matrix \(-[\int_0^1t^{i+j}\,dt]_{i,j=0}^m\). The
finite-dimensional implicit function theorem yields a continuously
differentiable coefficient map \(c(q,a)\). It is the minimizing
polynomial map by strict convexity and continuity already established.

The envelope derivative and its second derivative are

\[
\begin{aligned}
J_a&=q\int |R|^{q-2}R K_a^{[1]},\\
J_{aa}&=q(q-1)\int |R|^{q-2}
\bigl(K_a^{[1]}-P_{c_a}\bigr)K_a^{[1]}
+q\int |R|^{q-2}R K_a^{[2]},
\end{aligned}
\tag{6.14}
\]

where \(K_a^{[1]}=(a-t)_+^{m-1}/(m-1)!\) and
\(K_a^{[2]}=(a-t)_+^{m-2}/(m-2)!\). For \(m=2\), the last expression is
the indicator of \(t<a\), with its value at \(t=a\) immaterial. The
coefficient term in the first derivative vanishes by (6.13). The second
derivative exists and is continuous by (6.12), boundedness of these
kernels, and dominated convergence away from the moving interface. In
particular no claim that all residual roots are simple is needed.

At \(q=2\), the credited point-evaluation formula (4) gives

\[
J(2,a)=\frac{[a(1-a)]^{2n-1}}{(2n-1)((n-1)!)^2}.
\tag{6.15}
\]

It has a unique maximum at \(a=1/2\), with strictly negative second
derivative there. By continuity of (6.14), it follows that
\(J(q,\cdot)\) is strictly concave on one fixed neighborhood of \(1/2\),
for all \(q\) close enough to two. Uniform convergence on \([0,1]\)
forces every global maximum into that neighborhood. Reflection of the
original evaluation problem gives \(J(q,a)=J(q,1-a)\). Strict concavity
and this symmetry show that \(1/2\) is the unique global maximizing
point.

\subsection{6.4. Completion of the proof}\label{completion-of-the-proof}

Shrink \(\varepsilon_n\) so that both the variation-maximizer argument
and the evaluation-profile argument apply. We have established the
constant identity and coincidence of maximizing sets. Every height
maximizer must realize evaluation at its absolute maximum, and Section
6.3 forces that point to be \(1/2\). The norm-one positive evaluation
extremizer is unique by equality in Hölder in (6.4), or by strict
convexity of the unit ball of \(L^p\). Reflection preserves the norm and
evaluation at \(1/2\), so uniqueness makes this function midpoint
symmetric. Section 6.2 already gives its strict derivative signs and its
unique peak. This proves Theorem 5. \(\square\)

\subsection{6.5. The unresolved global
step}\label{the-unresolved-global-step}

\textbf{Compactness away from the endpoints of the exponent range.} The
argument in Lemma 6 gives the following more general statement. If
\(p_j\to p_0\in(1,\infty)\) and \(u_j\) are norm-one variation
maximizers, a subsequence converges in \(C^n\) to a norm-one variation
maximizer at \(p_0\), and \(V_{n,p_j}\to V_{n,p_0}\). To see this, keep
\(p_j\) in any fixed compact subinterval of \((1,\infty)\); the
exponents \(q_j-1\) then have a positive lower bound, so (6.6) still
gives a common Hölder modulus. The limit has norm one at \(p_0\). For
the lower bound on the limiting optimal value, test a maximizer at
\(p_0\); (6.6) makes its highest derivative bounded, so it is admissible
at all nearby exponents and its varying norms converge. The upper bound
follows by applying the \(p_0\) inequality to the limit. This proves the
statement without assuming uniqueness or symmetry at \(p_0\).

Thus a local continuation criterion at another exponent is available: if
all its normalized variation maximizers are scalar-sign copies of one
single-peak function with nonzero endpoint highest derivatives and a
nondegenerate peak, and its point-evaluation profile has a unique
midpoint maximum with a continuously persisting negative second
derivative, the same argument proves the conjecture in a neighborhood of
that exponent. These hypotheses must be established there; they cannot
be imported from the case \(p=2\).

The strongest compactness assertion here is \(C^n\) convergence of
\emph{variation maximizers}, established from their dual optimality
equation. It is not asserted for arbitrary normalized Sobolev sequences.
This distinction is necessary to preserve the derivative signs near the
clamped endpoints.

The point-profile calculation uses finite-dimensional best approximation
only after proving uniform control of possible multiple zeros in its
polynomial residual. Omitting (6.11)-\/-(6.12) would leave the two-sided
neighborhood in \(p\) unjustified, since \(q<2\) creates negative powers
at residual zeros.

This argument does not prove that the set of valid exponents is closed.
At a hypothetical first exponent where the result ceases to hold,
compactness of the old single-peak branch does not exclude a different
multiple-peak global maximizer with the same value. Nor does local
strict concavity of the point-evaluation profile exclude a remote
competing maximum at other exponents. These are substantive global
obstacles, not consequences of the implicit function theorem.

The remaining task is to exclude such competing branches or to find a
counterexample. Theorem 5 supplies a local result, not the full
conjecture. No numerical experiment is used in its proof.

\Needspace{8\baselineskip}\section{7. Remaining questions}\label{remaining-questions}

Theorems 1, 4 and 5 leave the full all-exponent form of {[}NS,
Conjecture 4.14{]} unresolved. For each order, Theorem 5 proves an open
interval about two, while Theorem 4 proves the constant identity and
coincidence at the measure endpoint. Symmetry and an explicit
classification of the maximizing knot configurations at p=1 remain open
here. Further continuation requires excluding competing maximizers or
finding a counterexample; the compactness argument alone does not do
that.

\raggedright\Needspace{8\baselineskip}\section{References}\label{references}

{[}ABR{]} S. Axler, P. Bourdon and W. Ramey, \emph{Harmonic Function
Theory}, second edition, Springer, 2001 (author PDF revision July 17,
2020). \href{https://axler.net/HFT.pdf}{Author text}, formula (1.15),
Theorem 1.17 and Chapter 5.

{[}BH{]} A. Burchard and H. Hajaiej, \emph{Rearrangement inequalities
for functionals with monotone integrands}, Journal of Functional
Analysis \textbf{233} (2006), 561-\/-582.
\href{https://doi.org/10.1016/j.jfa.2005.08.010}{DOI}.
\href{https://arxiv.org/abs/math/0506336}{Final revised author version}.

{[}dB{]} C. de Boor, \emph{Divided Differences}, Surveys in
Approximation Theory \textbf{1} (2005), 46-\/-69.
\href{https://arxiv.org/abs/math/0502036}{Author text}, formulas
(47)-\/-(48) and (52).

{[}GHW{]} R. J. Gardner, D. Hug and W. Weil, \emph{The
Orlicz-Brunn-Minkowski theory: A general framework, additions, and
inequalities}, Journal of Differential Geometry \textbf{97} (2014), no.
3. \href{https://doi.org/10.4310/jdg/1406033976}{DOI}.
\href{https://arxiv.org/abs/1301.5267}{Author text}.

{[}GS{]} T. A. Garmanova and I. A. Sheipak, \emph{Sharp Estimates of
High-Order Derivatives in Sobolev Spaces}, Moscow University Mathematics
Bulletin \textbf{79} (2024), 1-\/-10.
\href{https://doi.org/10.3103/S0027132224700013}{DOI}.

{[}HNOR{]} R. Hindov, S. Nitzan, J.-F. Olsen and E. Rydhe, \emph{A sharp
higher order Sobolev embedding}, Mathematika \textbf{71} (2025), e70012.
\href{https://doi.org/10.1112/mtk.70012}{DOI}.
\href{https://arxiv.org/abs/2411.10201}{Author preprint}.

{[}K{]} G. A. Kalyabin, \emph{Sharp Estimates for Derivatives of
Functions in the Sobolev Classes \(W_2^r(-1,1)\)}, Proceedings of the
Steklov Institute of Mathematics \textbf{269} (2010), 137-\/-142.
\href{https://doi.org/10.1134/S0081543810020112}{DOI}.

{[}NS{]} A. I. Nazarov and A. P. Shcheglova, \emph{A Survey of Results
on 1D Steklov Type Inequalities}, Proceedings of the Steklov Institute
of Mathematics \textbf{331} (2025), 134-\/-147; published online March
19, 2026. \href{https://doi.org/10.1134/S0081543825601509}{DOI}.
\href{https://arxiv.org/abs/2101.10752}{Latest author preprint}.

\end{document}
